\section{问题重述}

\subsection{问题背景}

干燥是决定中药材成品品质的关键工序之一\upcite{pin2023quality}，热风烘干则是其中最常见的实现方式，通常先经预热平衡使药材整体升温，再进入恒温干燥阶段持续脱水。烘干车间真正能够调控的只是烘房内的空气温度与含湿状况，而衡量成品是否合格的却是药材内部的含水率分布。这一层错位使得工艺参数的选取长期依赖经验：温度给低了、时间给短了，靠近轴心的部位仍然偏湿，储运中容易返潮霉变\upcite{dhulipalla2023storage}；温度给高了、时间拖长了，则表层过度失水甚至硬壳化，既浪费能耗又损伤有效成分\upcite{ai2022cistanche}。以反复试烘来标定参数，一炉动辄数十小时，成本高、周期长，难以覆盖不同品种与规格的组合；即便改用微波、真空等强化手段，工艺窗口仍需逐一实测标定\upcite{masztalerz2023adv,balbay2012liquorice}。

要摆脱这种试错模式，关键在于把烘房环境与药材内部状态之间的因果关系写成可计算的数学关系\upcite{vu2018transport}。药材内部并不存在可以直接读取的传感器读数，但热量由表面向内传导、水分由内部向表面迁移这两个过程本身遵循明确的守恒律，只要给定表面与空气的交换规律和材料的热质物性，内部的温度场与水分场原则上就被完全确定。困难在于三处：物性并非常数，密度、比热容、导热系数与水分扩散系数都随含水率甚至温度改变，方程因此是非线性的；合格判据针对的是药材各处而非平均状态，所求的是一个由场量首次全域达标所定义的时刻，而非某个显式表达式的取值；药材还会因失水而收缩，求解区域本身随时间变化。本题以一段圆柱形药材为研究对象，要求由烘房实测运行条件出发，逐步刻画预热阶段与全过程的温度与水分演化，并最终给出满足合格要求所需的烘干时长，为工艺参数的定量选取提供依据。

\subsection{问题的提出}

题目给出一段呈圆柱形的药材，其长度远大于半径，烘干开始时整体温度与含水率均匀，烘房的空气温度与水分浓度随时间变化的实测序列由附件给出。四个问题沿着“建立机理模型—扩展到全过程—引入合格判据—考虑几何收缩”的顺序递进，共用同一套传热传质机理，但在物性取法、求解区域与所求对象上依次变化。

问题一考察预热平衡阶段。在给定的初始温度与初始含水率下，以烘房实测的温度和水分浓度作为表面交换的驱动条件，采用该阶段适用的一组常物性参数与仅依赖含水率的扩散系数经验式，建立药材内部温度和水分浓度随时间与径向位置变化的数学模型，并按规定的时刻与径向位置给出半小时内的温度和含水率数值，同时输出时间与空间均加密的完整分布。这一问要确立后续各问共用的方程形式、边界条件与离散求解框架。

问题二把时间尺度从半小时延伸到整个烘干过程。烘干通常持续两到三天，预热平衡与恒温干燥两个阶段的材料参数并不相同，本问统一采用另一组经验公式，其中密度随含水率线性变化，比热容与导热系数随湿基含水率变化，扩散系数同时依赖含水率与绝对温度。模型结构继承问题一，但物性随解而变使方程转为非线性，需要相应的迭代求解策略；本问要求给出前三小时内的温度与含水率结果及完整分布，为后一问的停时判定提供全过程解。

问题三在问题二所得全过程解的基础上引入合格判据：药材各处的水分浓度都应低于给定阈值，据此确定烘干所需的时间。这里的关键是“各处”二字，它把问题从求解演化过程转变为确定一个首次达标时刻，判据应作用于药材内部含水率的最大值而非表面值或截面平均值。本问除给出烘干结束时间外，还要按规定的时间与位置间隔给出干燥过程中的水分浓度，并输出全程的完整水分场。

问题四考虑药材因水分流失而发生的尺寸变化。附件给出了烘干过程中药材半径随时间的实测序列，另有一组与收缩情形相对应的物性经验公式。此时求解区域的外边界随时间向轴线移动，前三问在固定区域上写出的方程不再直接适用，需要处理移动边界带来的几何变化，重新确定烘干时长并给出相应的水分浓度结果。本问与问题三的差别既来自物性组的更换，也来自几何收缩本身，两者的作用需要分离开来才能解释时长的变化。
